%% CV — Mahima Prajapati
%% Targeted for: Additive Manufacturing (AM) Internship — Edison Welding Institute / EWI (Buffalo, NY)
%% Compile with: pdflatex main_ewi.tex

\documentclass[11pt,a4paper]{moderncv}
\moderncvstyle{banking}
\moderncvcolor{blue}

\usepackage[utf8]{inputenc}
\usepackage{mathptmx}
\usepackage{hyperref}
\hypersetup{
    colorlinks=true,
    linkcolor=blue,
    filecolor=magenta,
    urlcolor=blue,
    pdftitle={Mahima Prajapati - CV},
    pdfpagemode=FullScreen,
}
\usepackage[scale=0.77]{geometry}
\usepackage{import}

% personal data
\name{Mahima}{Prajapati}
\address{Blacksburg, Virginia, USA}{}{}
\email{mahimap@vt.edu}
\extrainfo{\href{https://linkedin.com/in/mahimaprajapati}{LinkedIn}}

\begin{document}

\makecvtitle

% ============================================================
%     PROFILE STATEMENT
% ============================================================

\vspace{6pt}
\small{PhD candidate in Aerospace Engineering at Virginia Tech specializing in metal additive manufacturing simulation, with research directly focused on directed energy deposition (DED) and additive friction stir deposition (AFSD) --- processes central to EWI's advanced manufacturing R\&D capabilities. Develops Abaqus-based thermo-mechanical FEA models to predict deformations and distortions during metal AM processes, connecting computational simulation to real-world manufacturing outcomes. Published AIAA SciTech researcher (2026); second paper in preparation. Combines deep simulation expertise with a collaborative, research-oriented work style and a strong commitment to advancing manufacturing technology through rigorous, data-driven engineering.}

% ============================================================
%     CORE COMPETENCIES
% ============================================================

\section{Core Competencies}
\vspace{1pt}
\begin{itemize}

\item \textbf{Directed Energy Deposition (DED) \& Metal AM Simulation:} Abaqus (expert) --- thermo-mechanical FEA of laser wire DED and AFSD processes; predicting deformations, residual stresses, and distortions; validating models against experimental AM data.
\vspace{1pt}

\item \textbf{Materials Analysis \& Process Evaluation:} Evaluating material property behaviour in simulated AM components; thermo-mechanical coupling; process--structure--property relationships in metal AM.
\vspace{1pt}

\item \textbf{Computational Tools:} MATLAB --- numerical analysis, data processing, simulation scripting; ANSYS --- structural FEA validation.
\vspace{1pt}

\item \textbf{CAD \& Mechanical Design:} SolidWorks, CATIA --- 3D design and assembly; AutoCAD Electrical --- systems documentation; design for advanced manufacturing.
\vspace{1pt}

\item \textbf{Research Communication:} AIAA-published author (SciTech 2026); experienced technical communicator to 50+ students; journal paper in preparation; LaTeX documentation.

\end{itemize}

% ============================================================
%     EDUCATION
% ============================================================

\section{Education}
\vspace{1pt}
\begin{itemize}

\item{\cventry{Jan 2025--Present}{Doctor of Philosophy --- Aerospace \& Ocean Engineering}{Virginia Tech}{Blacksburg, VA}{}{\vspace{1pt}
Research: Abaqus-based thermo-mechanical FEA simulation of metal AM processes, with focus on laser wire DED and AFSD; predicting deformations and distortions; bridging computational modeling and physical manufacturing validation.
}}

\vspace{3pt}

\item{\cventry{Aug 2021--Dec 2024}{Master of Science --- Aerospace Engineering}{Virginia Tech}{Blacksburg, VA}{}{\vspace{1pt}
Thesis: ``Aeroelastic Analysis of a Uniform Cantilever Wing with Distributed Electric Propulsors with Rotary Inertia.'' Published at AIAA SciTech Forum 2026. \href{https://doi.org/10.2514/6.2026-1649}{DOI: 10.2514/6.2026-1649}
}}

\vspace{3pt}

\item{\cventry{2016--2020}{Bachelor of Engineering --- Mechanical Engineering}{M.H.\ Saboo Siddik College of Engineering}{Mumbai, India}{}{}}

\end{itemize}

% ============================================================
%     PROFESSIONAL EXPERIENCE
% ============================================================

\section{Professional Experience}
\vspace{3pt}
\begin{itemize}

\item{\cventry{Jan 2025--Present}{PhD Researcher --- Directed Energy Deposition \& Metal AM}{Virginia Tech, Kevin T. Crofton Dept. of Aerospace \& Ocean Engineering}{Blacksburg, VA}{}{\vspace{1pt}
\begin{itemize}
    \item Developing Abaqus thermo-mechanical FEA simulations of DED and AFSD processes; modeling heat transfer, residual stress, and part distortion across AM build sequences.
    \vspace{1pt}
    \item Collecting and analysing simulation and process data; evaluating material behaviour; iterating models to improve predictive accuracy against manufacturing outcomes.
    \vspace{1pt}
    \item Preparing manuscripts on DED/aeroelastic research for AIAA SciTech 2027 and a peer-reviewed journal; first paper published at AIAA SciTech 2026.
\end{itemize}}}

\vspace{3pt}

\item{\cventry{Aug 2021--Dec 2024}{Graduate Researcher \& Teaching Assistant --- Aeroelasticity}{Virginia Tech}{Blacksburg, VA}{}{\vspace{1pt}
\begin{itemize}
    \item Developed structural--aeroelastic simulation models for aircraft wings; conducted vibration and flutter analysis; published and presented findings at AIAA SciTech 2026 (\href{https://doi.org/10.2514/6.2026-1649}{DOI}).
    \vspace{1pt}
    \item Contributed to vibrational studies of the Broad Art Museum Monumental Staircase System with Dr.\ Mehdi Setareh; applied FEA methods to real-world structural problems.
    \vspace{1pt}
    \item Awarded \textbf{First Place --- Innovation \& Technology Advancement}, Virginia Tech GPSS (\$275 award); provided technical instruction to 50+ undergraduate students.
\end{itemize}}}


\end{itemize}

% ============================================================
%     PUBLICATIONS
% ============================================================

\section{Publications}
\vspace{1pt}
\begin{itemize}
\item Prajapati, M.~D., and Kapania, R.~K. (2026). ``Aeroelastic Analysis of a Uniform Cantilever Wing with Distributed Electric Propulsors with Rotary Inertia.'' \textit{AIAA SciTech Forum 2026}. \href{https://doi.org/10.2514/6.2026-1649}{DOI: 10.2514/6.2026-1649}
\item Prajapati, M.~D., and Kapania, R.~K. (2027). ``Effect of Thrust on Aeroelastic Analysis of DEP Wings with Rotary Inertia.'' \textit{AIAA SciTech Forum 2027}. (In preparation; journal paper forthcoming.)
\end{itemize}

% ============================================================
%     HONORS AND AWARDS
% ============================================================

\section{Honors and Awards}
\vspace{1pt}
\begin{itemize}
\item{\textbf{First Place --- Innovation \& Technology Advancement} -- Virginia Tech Graduate \& Professional Student Symposium (GPSS); \$275 monetary award.}
\item{\textbf{Tolson \& Nerney Fellowship} -- Kevin T. Crofton Dept.\ of Aerospace \& Ocean Engineering, Virginia Tech (2025); \$11,750.}
\end{itemize}

% ============================================================
%     REFERENCES
% ============================================================

\section{References}
\vspace{1pt}
\begin{itemize}
\item{Available upon request.}
\end{itemize}

\end{document}
